\clearpage
\fancyhead[L]{\small SOLAR PARALLAX.}
\SourceLeaf{528}
\begin{center}\small Göttingische gelehrte Anzeigen. Stück 17. pp. 161..166. January 29, 1810.\\[3pt]\rule{.72\linewidth}{.3pt}\end{center}

Mr \textsc{Lars Regner}, professor of astronomy at Uppsala and a member of the Royal Society of Sciences there, has sent in an essay entitled \emph{Supplementum ad historiam de Parallaxeos Solaris inventione, illustriss.\ scientiarum Societatis regiae Gottingensis dijudicationi subjectum}. Its substance is as follows. Mr. \textsc{Regner} first recounts that he had already been convinced in 1806 that the solar parallax accepted by astronomers was wrong, and that early the following year he had sent a short paper on the subject to the St Petersburg Academy of Sciences. Mr. \textsc{Fuss}, the Academy's secretary, had replied that the paper had been given to Mr. \textsc{Schubert} for assessment and that Schubert had found it erroneous and false. Mr. \textsc{Fuss} himself had no doubt that merely pointing out the source of the errors would suffice to convince Mr. \textsc{Regner} that his paper was wrong. He also offered to send the entire report itself. But Mr. \textsc{Regner} did not wait for it. Convinced that the criticism outlined to him was unjust, he had his paper printed in Uppsala with some explanations and then sent it to St Petersburg again. Mr. \textsc{Fuss} thereupon sent him Mr. \textsc{Schubert}'s report itself without further ado.

Mr. \textsc{Regner} next gives an extract from his printed paper (which has not otherwise come to our attention). Its main point is the observation that, if the Sun's parallax is assumed to lie between $8''5$ and $8''9$, the Sun's attraction on the Moon is considerably greater than the Earth's attraction on it. From this observation (which is indeed correct), Mr. \textsc{Regner} thinks he must conclude that this parallax cannot possibly be right. For the Moon, attracted more strongly by the Sun than by the Earth, would (on Mr. \textsc{Regner}'s view) have to leave the Earth and describe its own orbit around the Sun. Indeed, even if one assumed that the Moon was attracted by the Sun only as strongly as by the Earth, the two attractions would cancel each other completely at new moon. The Moon would move along the tangent and would thus immediately enter a region where the Sun's attraction predominates. Mr. \textsc{Regner} seeks to clarify this last claim with a figure in which he labels the Earth, the Moon and the Sun on a straight line by $T$, $b$, $S$ and then shows that the points on the perpendicular straight line erected at $b$ would recede from $T$ faster than from $S$. (How could Professor \textsc{Regner} forget that while the Moon moved along the tangent, the Earth would of course not remain at $T$, but would move along of its own accord by virtue of its tangential velocity, so the Moon's distance from it need not decrease at all? Likewise there is nothing contradictory in a \emph{greater} attraction toward the Sun. It follows \emph{only} that the Moon's orbit in absolute space must become concave toward the Sun at new moon. One can easily see that it really is so even by treating the orbit as an epicycloid and neglecting the inequalities in the motion of the Earth and the Moon.)

Mr. \textsc{Regner} now adds an extract from Mr. \textsc{Schubert}'s report. In \textsc{Schubert}'s judgment Mr. \textsc{Regner} must have written the paper either in great haste or in jest. For if he meant it seriously, it would follow that the author knew neither the method by which the solar parallax had been determined nor the \emph{first principles} of the three-body problem. \textsc{Schubert} supports the first claim by showing that one must assume an error of 27 minutes in the observations of the transit of Venus if the solar parallax were 21 seconds, which Mr. \SourceLeaf{529}\textsc{Regner}, on the view mentioned above, regards as the smallest admissible value. To substantiate the second claim, \textsc{Schubert} shows that Mr. \textsc{Regner} has confused the Moon's absolute motion with its motion relative to the Earth. In considering the latter, the Sun's attraction on the Moon is immaterial in itself. The sole consideration is how much it differs from the Sun's attraction on the Earth. This difference, however, is small compared with the Earth's attraction on the Moon. (Readers who are not altogether unfamiliar with the elements of physical astronomy will gladly excuse us from going further into \textsc{Schubert}'s explanation.)

Mr. \textsc{Regner} finally adds Animadversiones in relationem D\textsuperscript{nı} \textsc{Schubert} ad Scient. Imperial. Petropolitanam de Scripto parallaxin Solis tractante. We pass over his long complaints of the injustice by which he believes Mr. \textsc{Schubert} has hurt and insulted him, and merely give a brief account of how Mr. \textsc{Regner} seeks to rid himself of the charges. He first tries to cast doubt on the accuracy of the observations of the transits of Venus in 1761 and 1769, which \textsc{Lalande} had put at a few seconds of time or a few tenths of a second of arc. He maintains that such accuracy was simply unattainable with the telescopes used at the time and that \emph{for that reason} one must allow that the observational errors could well add up to 27 minutes of time, and declare the observations then made quite unfit for determining the solar parallax! (We cannot bring ourselves to add anything about such a way of defending oneself.)

Mr. \textsc{Regner} now thinks there is nothing left for him to do but refute Mr. \textsc{Schubert}'s second charge. He begins with a few commonplaces about the harm done by blind submission to authority and inherited errors in the sciences, for which he thinks this is just the place. He says it has ever since been believed, on \textsc{Newton}'s word alone, that at equal distances the effect of the Sun's attraction is equally great on the Earth and on the Moon. But he considers that proposition entirely false and puts the following in its place. When equal forces act on two bodies of unequal masses, the bodies acquire unequal velocities inversely proportional to their masses. (Why must we remind Professor \textsc{Regner} here that this proposition can be true only if we mean equal \emph{motive} force, as mechanics distinguishes it from accelerative force, and that the Sun exerts motive force on celestial bodies only insofar as it acts on their individual parts, thus exerting more motive force on the whole the more parts it attracts?) To defend this proposition Mr. \textsc{Regner} flatly denies the evidential force of the experiments with bodies falling under the bell jar of an air pump, partly because the times of fall are too short and partly because the air can never be \emph{entirely} removed. In Mr. \textsc{Regner}'s opinion the feather falls as fast as the ducat under the bell jar solely because the remaining air offers more resistance to the feather than to the gold coin. In a completely empty space the feather would fall faster by as much as its mass is smaller. He likewise thinks that the phenomenon of two spheres of equal volume and unequal masses, suspended by threads of equal length, making pendulum oscillations at the same rate (which, incidentally, is quite untrue insofar as the resistance of the air is still perceptible) is incontrovertible proof of his very proposition. For the smaller mass, he reasons, must be acted on by that much greater a force for it to overcome the same resistance as easily and fall as fast. Finally, Mr. \textsc{Regner} finds another proof of his proposition in the equilibrium of two bodies of unequal masses on unequal lever arms inversely proportional to those masses, because the virtual velocities here are in the inverse ratio of the masses. As Mr. \textsc{Regner} now takes all this to settle that \SourceLeaf{530}the effect of the Sun's attraction is greater on the Moon than on the Earth by as much as the former's mass is smaller than the latter's, he finds it easy to show that Mr. \textsc{Schubert}'s conclusions are inadmissible.

Neither the space nor the purpose of these pages allows us to go beyond this account and anticipate the conclusions that our readers will readily draw for themselves.
